% Fragment for the consolidated manuscript: TAUM's compact symmetry averaging and
% invariant compression (TAUM lines 1078-1187).
% Body only. Build with fragments/build_fragment.py symmetry_average.
\section{Averaging over a supplied symmetry}
\label{sec:symmetry-average}

A realization may come with a symmetry, a group that acts on the state space by unitary operators. The trap of Section~\ref{subsec:orbital-apparatus} is unchanged by rotations and by spatial parity. The pair Hamiltonian is unchanged when its two orbitals are exchanged. When the group is compact, averaging over it gives a projection onto the vectors that every element of the group leaves fixed. This section asks three questions. What is that projection? When does it respect the structures of the earlier sections, namely the sector symmetry, a preparation, and a Hamiltonian with its response? How is it related to the other projections in this manuscript?

The group and its action are supplied data. The average yields a projection and does not select a Hamiltonian. The symmetry is global. It limits what a symmetric preparation can reach, and nothing here makes it a local symmetry of a field theory on spacetime.

\subsection{The invariant projection}
\label{subsec:symmetry-projection}

Let $G$ be a compact Hausdorff group with its Haar probability measure $da$, which is invariant under left translation, right translation and inversion. Let $\mathcal K$ be a Hilbert space and let $\kappa:G\to\mathcal B(\mathcal K)$ be a strongly continuous representation by unitary operators. For $\xi,\eta\in\mathcal K$ the function $a\mapsto\langle\eta,\kappa(a)\xi\rangle$ is continuous and bounded by $\|\eta\|\,\|\xi\|$. Define a bounded operator $\Pi_G$ of norm at most one by
\begin{equation}
 \langle\eta,\Pi_G\xi\rangle=\int_G\langle\eta,\kappa(a)\xi\rangle\,da ,
 \label{eq:symmetry-average}
\end{equation}
and write the space of invariant vectors as
\begin{equation}
 \mathcal K^G=\{\xi\in\mathcal K:\kappa(a)\xi=\xi\text{ for every }a\in G\}.
 \label{eq:symmetry-invariant-space}
\end{equation}

\begin{proposition}[Averaging over a compact group]
\label{prop:symmetry-projection}
The operator $\Pi_G$ is the orthogonal projection onto $\mathcal K^G$, and
\begin{equation}
 \kappa(a)\,\Pi_G=\Pi_G\,\kappa(a)=\Pi_G\qquad\text{for every }a\in G .
 \label{eq:symmetry-absorption}
\end{equation}
It maps every closed subspace that is invariant under all $\kappa(a)$ into itself. If $G$ is finite, then
\begin{equation}
 \Pi_G=\frac1{|G|}\sum_{a\in G}\kappa(a).
 \label{eq:symmetry-finite-average}
\end{equation}
\end{proposition}
\begin{proof}
For $b\in G$, left invariance of the measure gives
\[
 \langle\eta,\kappa(b)\Pi_G\xi\rangle=\int_G\langle\kappa(b)^*\eta,\kappa(a)\xi\rangle\,da
   =\int_G\langle\eta,\kappa(ba)\xi\rangle\,da=\langle\eta,\Pi_G\xi\rangle ,
\]
so $\Ran\Pi_G\subseteq\mathcal K^G$. Right invariance gives $\Pi_G\kappa(b)=\Pi_G$ in the same way. If $\xi\in\mathcal K^G$ the integrand in Eq.~\eqref{eq:symmetry-average} is constant and $\Pi_G\xi=\xi$. Hence $\Pi_G^2=\Pi_G$ with range $\mathcal K^G$. Since $\kappa(a)^*=\kappa(a^{-1})$, invariance under inversion gives
\[
 \langle\Pi_G\eta,\xi\rangle=\int_G\langle\kappa(a)\eta,\xi\rangle\,da
   =\int_G\langle\eta,\kappa(a^{-1})\xi\rangle\,da=\langle\eta,\Pi_G\xi\rangle ,
\]
so $\Pi_G$ is self-adjoint. Let $\mathcal M$ be a closed invariant subspace, $\xi\in\mathcal M$ and $\eta\perp\mathcal M$. The integrand in Eq.~\eqref{eq:symmetry-average} vanishes, so $\Pi_G\xi$ is orthogonal to $\mathcal M^\perp$ and lies in $\mathcal M$. On a finite group the Haar probability measure gives each element the weight $1/|G|$.
\end{proof}

For a finite group of rotations of a cube, Eq.~\eqref{eq:symmetry-finite-average} is thus the exact projection onto the vectors on which that group acts trivially.

A smaller group fixes at least as many vectors. If $G'$ is a closed subgroup of $G$, then $\mathcal K^G\subseteq\mathcal K^{G'}$ and $\Pi_G\Pi_{G'}=\Pi_{G'}\Pi_G=\Pi_G$. Invariance under a finite subgroup and invariance under a continuous group are therefore different conditions. The 24 rotations of a cube act on the polynomials of degree four in three variables. They fix a two-dimensional space, spanned by $x^4+y^4+z^4$ and $x^2y^2+y^2z^2+z^2x^2$, whereas the full rotation group fixes only the multiples of $(x^2+y^2+z^2)^2$. A sequence of subgroups does recover the larger group in one case. If $G_1\subseteq G_2\subseteq\cdots$ are closed subgroups whose union is dense in $G$, the projections $\Pi_{G_n}$ decrease and converge strongly to $\Pi_G$, because the intersection of the spaces $\mathcal K^{G_n}$ consists of the vectors fixed by a dense subset of $G$, and these are fixed by all of $G$ by strong continuity. For finite groups that do not form such a sequence, such as the symmetry groups of finer and finer lattices, the convergence of the invariant projections, and of the Hamiltonians and responses compressed by them, is a separate problem that this manuscript leaves open.

\subsection{What the average respects}
\label{subsec:symmetry-compatibility}

\begin{proposition}[Sectors, preparations and Hamiltonians]
\label{prop:symmetry-compatibility}
\textup{(a)} Let $J$ be a symmetry of $\mathcal K$, $J=J^*=J^{-1}$, that commutes with every $\kappa(a)$. Then $\Pi_G$ commutes with $J$. On the direct sum $\mathcal H_\Psi\oplus\mathcal H_\Phi$ with $J=\operatorname{diag}(I_\Psi,-I_\Phi)$ this says that $\kappa=\kappa_\Psi\oplus\kappa_\Phi$ and
\begin{equation}
 \Pi_G=\Pi_G^\Psi\oplus\Pi_G^\Phi ,
 \label{eq:symmetry-sector-sum}
\end{equation}
the invariant projections of the two sectors. In the notation of Section~\ref{sec:neutral-graphs-phase}, $\Pi_G^\sharp=J\Pi_G^*J=\Pi_G$.

\textup{(b)} In the setting of \textup{(a)}, let $A:\operatorname{Dom}A\subseteq\mathcal H_\Psi\to\mathcal H_\Phi$ be a closed map that is covariant: $\kappa_\Psi(a)\operatorname{Dom}A\subseteq\operatorname{Dom}A$ and $A\kappa_\Psi(a)=\kappa_\Phi(a)A$ on $\operatorname{Dom}A$, for every $a\in G$. Then $\Pi_G^\Psi\operatorname{Dom}A\subseteq\operatorname{Dom}A$ and
\begin{equation}
 A\,\Pi_G^\Psi=\Pi_G^\Phi A\qquad\text{on }\operatorname{Dom}A .
 \label{eq:symmetry-covariant-map}
\end{equation}
Thus $A$ takes invariant input vectors to invariant output vectors. The adjoint $A^*$ is covariant in the same sense with the two sectors exchanged. The kernel of $A$ and the closure $D$ of its range are invariant, so the support projections of Section~\ref{sec:polar-matched-observables} commute with the representations and with the invariant projections.

\textup{(c)} Let $H$ be a self-adjoint operator in $\mathcal K$ whose resolvents commute with every $\kappa(a)$. Then $\mathcal K^G$ reduces $H$. Let $H^G$ be the restriction of $H$ to $\mathcal K^G$, and let $\Gamma:\mathcal E\to\mathcal K$ be a bounded port with $\Ran\Gamma\subseteq\mathcal K^G$. Then $\Hmin\subseteq\mathcal K^G$ and
\begin{equation}
 \Gamma^*(H-zI)^{-1}\Gamma=\Gamma^*(H^G-zI)^{-1}\Gamma ,\qquad z\in\C\setminus\R .
 \label{eq:symmetry-compressed-response}
\end{equation}
\end{proposition}
\begin{proof}
(a) For a bounded operator $T$ that commutes with every $\kappa(a)$,
\[
 \langle\eta,T\Pi_G\xi\rangle=\int_G\langle T^*\eta,\kappa(a)\xi\rangle\,da
   =\int_G\langle\eta,\kappa(a)T\xi\rangle\,da=\langle\eta,\Pi_GT\xi\rangle .
\]
Take $T=J$. An operator that commutes with $\operatorname{diag}(I_\Psi,-I_\Phi)$ is block diagonal, which gives the direct sums. The last identity follows from $\Pi_G^*=\Pi_G$ and $J^2=I$.

(b) The graph of $A$ is a closed subspace of $\mathcal H_\Psi\oplus\mathcal H_\Phi$, and covariance says that it is invariant under $\kappa_\Psi\oplus\kappa_\Phi$. By Proposition~\ref{prop:symmetry-projection} and Eq.~\eqref{eq:symmetry-sector-sum}, $\Pi_G^\Psi\oplus\Pi_G^\Phi$ maps the graph into itself, which is Eq.~\eqref{eq:symmetry-covariant-map}. For $\eta\in\operatorname{Dom}A^*$ and $\xi\in\operatorname{Dom}A$, covariance at $a^{-1}$ gives
\[
 \langle\kappa_\Phi(a)\eta,A\xi\rangle=\langle\eta,A\kappa_\Psi(a^{-1})\xi\rangle
   =\langle\kappa_\Psi(a)A^*\eta,\xi\rangle ,
\]
so $\kappa_\Phi(a)\eta\in\operatorname{Dom}A^*$ and $A^*\kappa_\Phi(a)\eta=\kappa_\Psi(a)A^*\eta$. The kernel of $A$ is closed and invariant, and so is $D$. The orthogonal complement of an invariant subspace of a unitary representation is invariant, so the projection onto it commutes with every $\kappa(a)$, and then with the invariant projection by the calculation in (a).

(c) By the calculation in (a), $\Pi_G$ commutes with every resolvent of $H$, so $\mathcal K^G$ reduces $H$, as in the proof of Proposition~\ref{prop:passive-two-spans}. Each vector $(H-zI)^{-1}\Gamma h$ lies in $\mathcal K^G$, and $\Hmin$ is the closed span of these vectors by that proposition. The resolvent of $H^G$ is the restriction of the resolvent of $H$, which gives Eq.~\eqref{eq:symmetry-compressed-response}.
\end{proof}

Part (c) is the sense in which a symmetry compresses a problem. A port that prepares only invariant vectors has the response of the compressed pair $(H^G,\Gamma)$. By Proposition~\ref{prop:preparation-split} every vector orthogonal to $\mathcal K^G$ is permanently invisible to it, and $H$ may be changed arbitrarily on $(\mathcal K^G)^\perp$, also in ways that break the symmetry, without changing that response. The rotations $\exp(-i\varphi L/\hbar)$ of the trap form a circle group. Spatial parity in the plane is the rotation by $\pi$, and the even-parity space used for the shell in Proposition~\ref{prop:shell-accuracy} is the invariant space of the subgroup of two elements that it generates. The invariant projection of the whole circle group is the spectral projection of $L$ at zero, and the eigenspace $L=\ell$ of Proposition~\ref{prop:passive-opposite-sectors} is the range of the twisted average $(2\pi)^{-1}\int_0^{2\pi}e^{i\ell\varphi/\hbar}\exp(-i\varphi L/\hbar)\,d\varphi$.

The projection $\Pi_G$ is a different object from the projection onto the input sector, from the support projections of a map, and from the projection onto a prepared or moving subspace in Section~\ref{sec:moving-subspace}. Its construction does not identify it with any of them. Such an identification is additional data of a realization and has to be checked in each case, as the following example shows.

\subsection{Two symmetries of the pair model}
\label{subsec:symmetry-pair}

On the two-particle space $\mathcal F_2$ of the pair model, let $\sigma$ be the unitary that exchanges the two orbitals, $c_i\leftrightarrow d_i$ for $i=1,2$, and fixes the vacuum. It commutes with $H_s$ for every $s$. It fixes $g$ and $m$, and it reverses $e$ and every vector of $W$. The group $\{I,\sigma\}$ therefore has the invariant projection
\begin{equation}
 \Pi_\sigma=\tfrac12(I+\sigma)=\ket{g}\bra{g}+\ket{m}\bra{m}=P_{\pi/4},
 \label{eq:symmetry-pair-exchange}
\end{equation}
where $P_\theta=A_\theta A_\theta^*$ projects onto the prepared plane and $P_{\pi/4}$ onto the plane $D_{\pi/4}$ of Eq.~\eqref{eq:pair-two-setting-ranges}. At that setting the invariant projection is the prepared projection, and part (c) of Proposition~\ref{prop:symmetry-compatibility} confines its orbit to the invariant plane. This agrees with Section~\ref{subsec:moving-pair}, where the plane prepared at this setting does not change with the waiting time. At $\theta=0$ the prepared plane $\Span\{g,e\}$ commutes with $\Pi_\sigma$ and differs from it. For $0<\theta<\pi/4$ the excitation $b=u_\theta e+v_\theta m$ mixes the two eigenspaces of $\sigma$, and $P_\theta$ does not commute with $\Pi_\sigma$.

The preparation itself comes from a second group. Let $\kappa(\varphi)$ rotate the orbitals, $c_i^*\mapsto\cos\varphi\,c_i^*+\sin\varphi\,d_i^*$ and $d_i^*\mapsto-\sin\varphi\,c_i^*+\cos\varphi\,d_i^*$. Then $A_\theta=\kappa(\theta)A_0$ and $\kappa(\theta)e=b$. On $\mathcal F_2$ the rotation fixes $g$ and every vector of $W$, and it turns the pair $(e,m)$ through the angle $2\varphi$. Its invariant projection is
\begin{equation}
 \Pi_{\rm rot}=\frac1{2\pi}\int_0^{2\pi}\kappa(\varphi)\,d\varphi=\ket{g}\bra{g}+\Pi_W ,
 \label{eq:symmetry-pair-rotation}
\end{equation}
of rank four. The prepared excitation is orthogonal to the invariant space at every setting, $\Pi_{\rm rot}P_\theta=\ket{g}\bra{g}$. The rotations commute with $H_s$ only at $s=j$, where $H_s=-2j\ket{g}\bra{g}$. For $s\neq j$ part (c) does not apply to them, which is the statement of Section~\ref{sec:preparation-effective-dynamics} that the Hamiltonian breaks the preparation rotation.

The two groups illustrate the difference between finite and continuous invariance. The rotation by $\pi$ acts as the identity on $\mathcal F_2$, so averaging over the subgroup $\{\kappa(0),\kappa(\pi)\}$ gives the identity, of rank six. The cyclic subgroup of order $n\geq3$ already gives $\Pi_{\rm rot}$, because it resolves the two units of rotational charge carried by the pair $(e,m)$. The exchange is not a rotation, since $\sigma\kappa(\varphi)\sigma=\kappa(-\varphi)$. Together they generate the orthogonal group of the orbital pair, whose invariant projection is $\Pi_\sigma\Pi_{\rm rot}=\ket{g}\bra{g}$. The ground $g$ is the only two-particle state fixed by every orthogonal transformation of the two orbitals.
